En el fragmento evaluado, el aumento de la atenuación del eco produjo una disminución del SNR, de $42.581$ dB con un $5\,\char37{}$ de atenuación a $28.605$ dB con un $25\,\char37{}$. Por tanto, existe un compromiso entre la intensidad del eco insertado y la alteración de la señal. La selección de una atenuación adecuada requiere comprobar también la recuperación de los bits. Los resultados actuales corresponden a un único fragmento musical y a un mensaje de seis bits; no permiten generalizar el comportamiento a otros géneros ni a una ocupación completa de las ventanas.
